% Options for packages loaded elsewhere
\PassOptionsToPackage{unicode}{hyperref}
\PassOptionsToPackage{hyphens}{url}
\PassOptionsToPackage{dvipsnames,svgnames,x11names}{xcolor}
%
\documentclass[
  11pt,
]{article}
\usepackage{amsmath,amssymb}
\usepackage{iftex}
\ifPDFTeX
  \usepackage[T1]{fontenc}
  \usepackage[utf8]{inputenc}
  \usepackage{textcomp} % provide euro and other symbols
\else % if luatex or xetex
  \usepackage{unicode-math} % this also loads fontspec
  \defaultfontfeatures{Scale=MatchLowercase}
  \defaultfontfeatures[\rmfamily]{Ligatures=TeX,Scale=1}
\fi
\usepackage{lmodern}
\ifPDFTeX\else
  % xetex/luatex font selection
\fi
% Use upquote if available, for straight quotes in verbatim environments
\IfFileExists{upquote.sty}{\usepackage{upquote}}{}
\IfFileExists{microtype.sty}{% use microtype if available
  \usepackage[]{microtype}
  \UseMicrotypeSet[protrusion]{basicmath} % disable protrusion for tt fonts
}{}
\makeatletter
\@ifundefined{KOMAClassName}{% if non-KOMA class
  \IfFileExists{parskip.sty}{%
    \usepackage{parskip}
  }{% else
    \setlength{\parindent}{0pt}
    \setlength{\parskip}{6pt plus 2pt minus 1pt}}
}{% if KOMA class
  \KOMAoptions{parskip=half}}
\makeatother
\usepackage{xcolor}
\usepackage[margin=1in]{geometry}
\usepackage{longtable,booktabs,array}
\usepackage{calc} % for calculating minipage widths
% Correct order of tables after \paragraph or \subparagraph
\usepackage{etoolbox}
\makeatletter
\patchcmd\longtable{\par}{\if@noskipsec\mbox{}\fi\par}{}{}
\makeatother
% Allow footnotes in longtable head/foot
\IfFileExists{footnotehyper.sty}{\usepackage{footnotehyper}}{\usepackage{footnote}}
\makesavenoteenv{longtable}
\setlength{\emergencystretch}{3em} % prevent overfull lines
\providecommand{\tightlist}{%
  \setlength{\itemsep}{0pt}\setlength{\parskip}{0pt}}
\setcounter{secnumdepth}{-\maxdimen} % remove section numbering
\ifLuaTeX
  \usepackage{selnolig}  % disable illegal ligatures
\fi
\usepackage{bookmark}
\IfFileExists{xurl.sty}{\usepackage{xurl}}{} % add URL line breaks if available
\urlstyle{same}
\hypersetup{
  pdftitle={P16.3: positive residue characteristic, source audit and a partial theorem},
  pdfauthor={Research draft},
  colorlinks=true,
  linkcolor={blue},
  filecolor={Maroon},
  citecolor={Blue},
  urlcolor={Blue},
  pdfcreator={LaTeX via pandoc}}

\title{P16.3: positive residue characteristic, source audit and a
partial theorem}
\author{Research draft}
\date{4 October 2026}

\begin{document}
\maketitle

\textbf{Status: partial; the original nontrivially valued
positive-residue-characteristic atom is not solved.} This document
proves a precisely delimited application of published results and two
elementary normalization/variational passages. It does not claim a new
extension of the non-Archimedean finite-energy equation theorem. In
particular, the higher-dimensional nontrivially valued case remains an
obligation even for smooth projective varieties. All source references
below refer to the ten supplied, hash-verified PDFs, not to later
versions inferred from an abstract.

\section{1. The target and the valuation
distinction}\label{the-target-and-the-valuation-distinction}

The residual P16.3 target retains a complete \textbf{nontrivially
discretely valued} non-Archimedean field \(K\), a smooth projective
variety \(X/K\), dimension \(n\geq 1\), and an ample line bundle \(L\).
Write \(V=(L^n)>0\). The restriction to remove is
\(\operatorname{char}(\widetilde K)=0\), where \(\widetilde K\) is the
residue field. Thus both equal characteristic \((p,p)\) and mixed
characteristic \((0,p)\) must be addressed. Field characteristic zero
alone does not meet the original residue-characteristic hypothesis. For
uniqueness by one constant one works on a connected variety;
disconnected components instead permit independent constants. The
unresolved case already contains geometrically integral smooth
projective varieties, so introducing that restriction would not close
the gap.

The required equation is

\[\operatorname{MA}^{\rm raw}_{\rho}(u)=\mu,\qquad \mu\geq 0,\quad \mu(X^{\rm an})=V,\]

for an actual psh potential \(u\) of finite energy relative to a
specified continuous semipositive reference metric \(\rho\) on \(L\).
The positivity datum remains the ample \(L\), and its raw mass remains
\(V\). Measures must satisfy the dual-energy condition specified below.
No support restriction to a skeleton, toric invariance, or replacement
inverse measure is authorized. There must be an
approximation-independent extension of the genuine model Monge--Ampere
measure, followed by existence and uniqueness modulo constants on the
stated finite-energy space.

If \(k\) has the trivial absolute value, then \(k^\circ=k\),
\(k^{\circ\circ}=0\), and \(\widetilde k=k\). Thus a trivially valued
field of characteristic \(p\) has residue characteristic \(p\), but it
is not a nontrivially valued field. Neither \(\overline{\mathbb F_p}\)
with trivial valuation nor its test configurations give a proof for
\(\mathbb F_p((t))\), much less \(\mathbb Q_p\). Test configurations are
auxiliary models over \(k[t]\), not an identification of an arbitrary
nontrivially valued variety with a trivially valued one.

The following are the conclusions warranted by this investigation:

\begin{longtable}[]{@{}
  >{\raggedright\arraybackslash}p{(\columnwidth - 4\tabcolsep) * \real{0.3333}}
  >{\raggedright\arraybackslash}p{(\columnwidth - 4\tabcolsep) * \real{0.3333}}
  >{\raggedright\arraybackslash}p{(\columnwidth - 4\tabcolsep) * \real{0.3333}}@{}}
\toprule\noalign{}
\begin{minipage}[b]{\linewidth}\raggedright
Domain
\end{minipage} & \begin{minipage}[b]{\linewidth}\raggedright
Established conclusion used here
\end{minipage} & \begin{minipage}[b]{\linewidth}\raggedright
Limitation
\end{minipage} \\
\midrule\noalign{}
\endhead
\bottomrule\noalign{}
\endlastfoot
Algebraically closed trivially valued \(k\), characteristic \(p>0\),
smooth connected projective \(X\), ample \(L\), every dimension &
Genuine \(\mathcal E^1\) operator, strong continuity, comparison, and
injectivity modulo constants & Surjectivity in arbitrary dimension is
not supplied unconditionally \\
The preceding domain with \(1\leq n\leq2\) & Every positive Radon
measure of mass \(V\) and finite dual energy has a unique finite-energy
solution modulo constants & A trivial-valuation endpoint, including
surfaces, not P16.3's nontrivial axis \\
Smooth projective \(X\) over any complete non-Archimedean \(K\),
including positive residue characteristic & Genuine MA for continuous
semipositive metrics; test-function dual energy, orthogonality in the
synthetic sense, and strong density of test MA measures &
Density/completeness on the measure side is not existence of an
unbounded potential \\
Discrete equal characteristic \(p\), function-field descent of
\((X,L)\), and assumed resolution in dimension \(n+d\) & A
positive-characteristic envelope result is reported in BE, Section 7.30
& Additional descent and resolution hypotheses; no independent proof of
the full finite-energy equation here \\
Nontrivial discrete valuation, residue characteristic \(p\),
arbitrary-dimensional smooth projective \(X\) and ample \(L\), with no
such extra assumptions & Full P16.3 theorem not established in this
draft & Concrete replacement obligations in Section 7 \\
\end{longtable}

These limitations describe what is established here and in the checked
source statements. They are not a proof that no future or uncached
theorem exists, and they are not a counterexample to P16.3.

\section{2. Exact versions and the original proof
dependencies}\label{exact-versions-and-the-original-proof-dependencies}

\textbf{BFJ-PSH} denotes Boucksom--Favre--Jonsson, \emph{Singular
semipositive metrics in non-Archimedean geometry}, arXiv:1201.0187v2.
Its PDF banner is 21 January 2014; the internal manuscript date says 22
October 2018. Both are retained in provenance. It assumes complete
discrete valuation, residue characteristic zero, smooth projective \(X\)
and ample \(L\). In the introduction and Sections 3, 5--8 it uses the
following inputs.

\begin{enumerate}
\def\labelenumi{\arabic{enumi}.}
\tightlist
\item
  Cofinal regular SNC models: every model is dominated by a regular
  model with special fiber of SNC support. The inverse-limit skeleton
  description and the retractions enter its original definition of psh
  and approximation by retractions. Generic-fiber smoothness does not
  imply regularity or SNC support for the total model.
\item
  Numerical/model positivity: Theorem 5.11 includes multiplier ideals in
  its positivity argument. Resolving the model does not replace this use
  of ideals.
\item
  Compactness: Theorem 6.1 supplies equicontinuity on dual complexes,
  using toroidal modifications and intersection estimates. Section 7
  passes to compactness modulo constants. This compactness is
  compactness of the psh cone, not merely compactness of the underlying
  analytic space.
\item
  Envelopes and regularization: Theorems 8.3, 8.5 and 8.7 use graded
  base ideals, asymptotic multiplier ideals, subadditivity and uniform
  global generation. Appendix B proves Kodaira/Kawamata--Viehweg/Nadel
  vanishing and ideal properties on the regular/SNC models in the
  standing characteristic-zero regime. Deleting that regime does not
  validate these arguments.
\end{enumerate}

\textbf{BFJ-MA-v1} is the actual saved arXiv:1201.0188v1, banner 30
December 2011. Theorem A additionally requires that \(X\) be
\emph{algebraizable}: \(X=Y_K\) for a smooth projective variety \(Y\)
over a one-variable function field \(F\) admitting \(K\) as a
completion. This differs from saying that \(X\) is an algebraic
\(K\)-variety. The theorem does \textbf{not} require that \(L\) descend
to \(F\). It requires \(\mu\) to be supported on the dual complex of an
SNC model, and requires mass \((L^n)\). Its conclusion is a continuous
semipositive solution, unique up to multiplication of the metric by a
positive constant.

The v1 proof then depends on BFJ-PSH compactness and regularization to
construct bounded MA, locality, truncation measures and the
\(\mathcal E^1\) operator. Sections 3--6 are the
operator/locality/energy part; Section 7 differentiates energy composed
with the envelope, and Section 8 obtains solutions. Appendix A, Theorem
A.4, uses the function-field algebraizability hypothesis to compactify a
model into a projective variety over the residue field and apply the
orthogonality of big-line-bundle decompositions. The introduction
explicitly identifies this role. SNC support is also used for its
continuity estimates; it cannot be read as an unrestricted finite-energy
measure theorem just because the proof introduces \(\mathcal E^1\).

\textbf{BE} is Boucksom--Eriksson, \emph{Spaces of norms, determinant of
cohomology and Fekete points in non-Archimedean geometry},
arXiv:1805.01016v3, banner 8 January 2021. It is not a differently
titled pluripotential paper. For geometrically reduced projective \(X\)
over any complete valued field and semiample \(\mathbb Q\)-line bundle
\(L\), Theorem 9.15 identifies relative volume with relative MA energy
for two continuous psh metrics. Smooth \(X\) satisfies geometric
reducedness, and ample \(L\) is semiample. Section 7.30 still labels
unrestricted envelope continuity a conjecture. Its reported
characteristic-\(p\) discrete case assumes descent of \textbf{\((X,L)\)}
to a function field of transcendence degree \(d\) and resolution in
dimension \(n+d\), replacing multiplier ideals with test ideals. The
same restrictions recur in Section 10.9's differentiability discussion.
The cited GJKM paper is not among the ten allowed inputs; its exact
stronger resolution formulation is consequently not asserted or newly
proved here. The reported statement cannot cover mixed characteristic,
nor all non-descending pairs.

\textbf{BJ-triv} is Boucksom--Jonsson, \emph{Global pluripotential
theory over a trivially valued field}, arXiv:1801.08229v4, banner 22
March 2022. Its main-body ground field is algebraically closed, of
arbitrary characteristic, with trivial valuation. Theorems 7.18, 8.5,
10.4, 12.8 and 12.9 have distinct roles described below.

\textbf{BJ-synth} is Boucksom--Jonsson, \emph{Measures of finite energy
in pluripotential theory: a synthetic approach}, arXiv:2307.01697v1, 6
July 2023. The introduction explicitly says it does not investigate
potentials of finite energy. Its theorem about completion of measures is
therefore not a theorem identifying that completion with
\(\mathcal E^1/\mathbb R\).

\textbf{PS-hybrid}, arXiv:2209.04879v2, contains the useful distinction
in Remark 4.4: defining \(\mathcal E^1\) over any non-Archimedean field
does not generally give its MA extension. Its Section 1.2 expressly
assumes \emph{equicharacteristic zero}. The shorter words
``characteristic zero'' in Remark 4.4 cannot enlarge this standing
setting to mixed characteristic.

\section{3. Partial theorem over a trivially valued
field}\label{partial-theorem-over-a-trivially-valued-field}

\textbf{Theorem (published endpoint in raw normalization).} Let \(k\) be
algebraically closed of characteristic \(p>0\), equipped with the
trivial absolute value. Let \(X/k\) be smooth, connected and projective,
of dimension \(n\geq1\), and let \(L\) be ample. Fix the canonical
trivial reference metric \(\rho\), and put \(\omega=c_1(L)\),
\(V=(L^n)\). There is a genuine, approximation-independent operator

\[\operatorname{MA}^{\rm raw}_{\rho}:\mathcal E^1(L)\longrightarrow\{\text{positive Radon measures of mass }V\}.\]

It extends the intersection-theoretic model MA measure, is continuous
along decreasing nets, and its normalized version is continuous for the
strong topologies. Its image consists of finite-dual-energy measures; it
is injective modulo constants, with dense image in that measure space.
If \(1\leq n\leq2\), its image is the entire finite-dual-energy measure
space. Equivalently, for every positive Radon \(\mu\) with mass \(V\)
and \(\mathcal E_{\rho}^{*,\rm raw}(\mu)<\infty\), there is a solution
\(u\in\mathcal E^1(L)\), unique modulo addition of a real constant.
There is no restriction on measure support. No assertion that \(u\) is
continuous is made. For \(n\geq3\), the same surjectivity follows
\textbf{if} \(L\) has the envelope property, but that extra hypothesis
is not proved here.

This is a source theorem application, not a newly proved removal of the
nontrivial residue-characteristic restriction.

\subsection{3.1 Positivity, reference, functions and exact measure
energy}\label{positivity-reference-functions-and-exact-measure-energy}

The canonical trivial metric is the reference used in BJ-triv's
Fubini--Study functions. Define \(\mathcal H(L)\) using

\[u=\frac1m\max_i(\log|s_i|+\lambda_i),\qquad s_i\in H^0(X,mL),\quad \lambda_i\in\mathbb Q,\]

where the sections have no common zero. These are continuous
semipositive metrics relative to \(\rho\) and correspond to ample
integrally closed test configurations. Define \(\operatorname{PSH}(L)\)
as the usc functions, not identically \(-\infty\), obtained as
decreasing nets from \(\mathcal H(L)\). They are stable under constants
and maxima and are determined by their finite values on divisorial
valuations. The weak topology is pointwise convergence on these
valuations. The trivial valuation \(v_{\rm triv}\) satisfies
\(\sup u=u(v_{\rm triv})\). The canonical reference has normalized MA
equal to \(\delta_{v_{\rm triv}}\).

For continuous semipositive \(u\) put

\[\mathcal E_{\rho}^{\rm raw}(u)=\frac1{n+1}\sum_{j=0}^n\int_{X^{\rm an}}u\, (\omega+dd^c u)^j\wedge\omega^{n-j},\qquad E(u)=V^{-1}\mathcal E_{\rho}^{\rm raw}(u).\]

Here the mixed products are the genuine intersection-theoretic measures,
and the formula agrees with the compactified-test-configuration energy.
Extend energy monotonically by

\[E(u)=\inf\{E(v):v\in\mathcal H(L),\ v\geq u\},\qquad \mathcal E^1(L)=\{u\in\operatorname{PSH}(L):E(u)>-\infty\}.\]

BJ-triv Proposition 7.7 supplies concavity, upper semicontinuity,
decreasing-net continuity and \(E(u+c)=E(u)+c\). The strong topology is
the refinement of the weak topology making \(E\) continuous. On the
measure side it is weak convergence together with convergence of dual
energy.

The exact raw condition is

\[\mathcal E_{\rho}^{*,\rm raw}(\mu)=\sup_{u\in\mathcal H(L)}\left\{\mathcal E_{\rho}^{\rm raw}(u)-\int u\,d\mu\right\}<\infty,\qquad \mu(X^{\rm an})=V.\]

Writing \(\nu=\mu/V\) is simply bookkeeping for a measure already of the
correct mass; it is not normalization of an arbitrary restriction of
\(\mu\). Exactly

\[\mathcal E_{\rho}^{*,\rm raw}(\mu)=V E^\vee(\nu),\qquad E^\vee(\nu)=\sup_{u\in\mathcal E^1}\left(E(u)-\int u\,d\nu\right)=\sup_{u\in\mathcal H}\left(E(u)-\int u\,d\nu\right).\]

\textbf{Proof of the raw passage and integrability.} The displayed
equality follows by factoring \(V\) from each test expression. If the
test supremum is finite and \(u_j\in\mathcal H\) decreases to
\(u\in\mathcal E^1\), then \(E(u_j)\to E(u)>-\infty\) and

\[\int u_j\,d\nu\geq E(u_j)-E^\vee(\nu).\]

The usc decreasing-net integration theorem for a positive Radon measure
on the compact \(X^{\rm an}\) (BJ-triv Lemma 7.17) gives
\(\int u_j\,d\nu\to\int u\,d\nu\). A uniform upper bound for \(u_j\) and
the displayed lower bound show that \(u\in L^1(\nu)\). Each
finite-energy test expression is therefore the limit of its model
expressions and is bounded by the test supremum. The reverse supremum
inequality follows from \(\mathcal H\subset\mathcal E^1\). This proves
the equality, with no assumed solution to the equation. If an arbitrary
positive measure has mass \(M\ne V\), the constant tests give
\(c(V-M)\), whose supremum over \(c\in\mathbb R\) is infinite. Thus
finite raw dual energy itself forces exactly the required mass. In the
trivial setting BJ-triv Theorem 9.20 also gives the equivalent condition
\(\mathcal E^1\subset L^1(\nu)\).

Changing the reference to a continuous semipositive metric \(\rho+r\)
changes the raw dual functional by the finite quantity
\(\int r\,d\mu-\mathcal E_{\rho}^{\rm raw}(r)\). Indeed, the
relative-energy cocycle yields
\(\mathcal E_{\rho+r}^{\rm raw}(u-r)=\mathcal E_{\rho}^{\rm raw}(u)-\mathcal E_{\rho}^{\rm raw}(r)\),
and substitution in the supremum proves the assertion. This passage uses
the published relative-energy identity on the already established
endpoint; it asserts no unbounded descent over another field.

\subsection{3.2 Construction of the genuine operator and
independence}\label{construction-of-the-genuine-operator-and-independence}

Use the multilinear finite-energy pairing of BJ-triv Theorem 7.14. Its
existence, finiteness and monotone continuity are published inputs,
proved there by a monotone extension of the PL intersection pairing and
mixed-energy estimates. For \(u\in\mathcal E^1\) define on real PL test
functions \(f\)

\[T_u(f)=(0,f)\cdot(\omega,u)^n.\]

Lemma 7.15 and Theorem 7.18 make \(T_u\) a positive linear functional.
Constants satisfy \(T_u(1)=V\) by the constant-shift identity of the
pairing. Thus \(|T_u(f)|\leq V\|f\|_\infty\). Real PL functions are
uniformly dense in \(C^0(X^{\rm an})\) by Theorem 2.2. Consequently
\(T_u\) extends uniquely to a positive bounded functional on \(C^0\).
Riesz representation gives its unique positive Radon measure, which we
define as \(\operatorname{MA}_{\rho}^{\rm raw}(u)\).

If \(u_j\) decreases to \(u\) through model functions, monotone
continuity of the pairing gives \(T_{u_j}(f)\to T_u(f)\) for each PL
\(f\). The common bound \(V\|f\|_\infty\) extends convergence to every
continuous \(f\). Therefore the weak limit is independent of the chosen
net. On a model/test configuration with components \(F\) in its special
fiber, a model metric \(u\) has

\[\operatorname{MA}_{\rho}^{\rm raw}(u)=\sum_F b_F\bigl((\omega_{\mathcal X}+D)|_F\bigr)^n\delta_{v_F},\qquad b_F=\operatorname{ord}_F(\mathcal X_0),\]

as in Proposition 7.19(ii); hence this operator extends the genuine
model MA. Proposition 7.19(i) gives raw total mass \(V\). Theorem A
gives strong continuity. Neither a norm-derived inverse measure nor an
arbitrary assignment of an atomic distribution has been substituted. No
resolution is needed in the construction in BJ-triv Section 7;
integrally closed test configurations and the pairing replace the
original SNC-dependent construction in this trivial setting.

\subsection{3.3 Envelopes and compactness: the actual characteristic
replacement}\label{envelopes-and-compactness-the-actual-characteristic-replacement}

For a continuous obstacle \(h\) put
\(P(h)=\sup\{u\in\operatorname{PSH}(L):u\leq h\}\). The envelope
property means that usc regularization of the supremum of every
bounded-above psh family remains psh. BJ-triv Lemma 5.17 identifies it
with continuity of \(P(h)\) for continuous \(h\); Theorem 5.11
identifies it with weak compactness of the normalized psh space.

For \(n\leq2\), Theorem 5.20 proves this property in positive
characteristic. Its proof first reduces to a PL obstacle represented by
\((\mathcal X,\mathcal L)\) and writes

\[P(\varphi_{\mathcal L})-\varphi_{\mathcal L}=\lim_m m^{-1}\varphi_{\mathfrak a_m},\]

where \(\mathfrak a_m\) is the base ideal of \(m\mathcal L\) (Lemma
5.22). Normalized blowups of these ideals give the nef approximants. The
proof resolves the test configuration to a smooth total space with SNC
central support. The relevant dimension is
\textbf{\(\dim\mathcal X=n+1\leq3\)}, not just the dimension of \(X\).
The supplied PDF page 62 visually distinguishes the roman \(X\) from the
calligraphic \(\mathcal X\), which the cached OCR flattens.

In characteristic zero the proof uses asymptotic multiplier ideals; in
positive characteristic it uses test ideals and their
subadditivity/global generation, as expressly stated on PDF page 63. The
resulting estimate has the form

\[\varphi_m\leq P(\varphi_{\mathcal L})\leq\left(1+\frac{m_0}{m}\right)\varphi_{m+m_0}-\frac{m_0}{m}\varphi_{\mathcal L}+\frac C m.\]

It forces uniform convergence and hence envelope continuity. This
paragraph reconstructs the source's replacement mechanism; the algebraic
resolution and test-ideal results used in its published Theorem 5.20 are
background inputs, not reproved results of this draft. Crucially, their
available low-dimensional use is why the unconditional
positive-characteristic theorem stated here has \(n\leq2\). Assuming
resolutions in every necessary dimension would add a hypothesis. No
unrestricted higher-dimensional resolution theorem is claimed.

Once the envelope property is known, normalized psh compactness can be
seen directly. For a normalized net, divisorial values are uniformly
bounded by the source's evaluation estimates. Tychonoff gives a subnet
converging at every divisorial point. The usc tail suprema are psh by
the envelope property and decrease to the usc potential having those
limits; nontriviality follows from finite divisorial values. The
divisorial determination theorem identifies that limit uniquely. This is
the proof mechanism of Theorem 5.11 and avoids invoking cofinal SNC
skeletons as the compactness argument.

\subsection{3.4 Energy differentiation, comparison and
uniqueness}\label{energy-differentiation-comparison-and-uniqueness}

BJ-triv Theorem 8.5 and Corollary 8.6 prove, for \(u\in\mathcal E^1\),
PL \(f\) and arbitrary continuous \(g\),

\[E^\downarrow(u+tf)=E(u)+t\int f\,\operatorname{MA}(u)+O(t^2),\qquad \left.\frac d{dt}\right|_{t=0}E^\downarrow(u+tg)=\int g\,\operatorname{MA}(u).\]

Here \(\operatorname{MA}=V^{-1}\operatorname{MA}^{\rm raw}\) and
\(E^\downarrow\) is the decreasing extension of obstacle energy. When
the envelope property holds, \(E^\downarrow(u+tg)=E(P(u+tg))\). The
proof uses relative volumes and a uniform differentiability estimate,
followed by monotone approximation; it is not an assertion that an
arbitrary perturbation is psh. Multiplying by \(V\) gives the raw
derivative \(\int g\,\operatorname{MA}^{\rm raw}(u)\). The source's
proof cites BE and the relative-volume differentiability work; these are
published background, with BE's geometric reducedness satisfied here.
The differentiated statement used is the fully stated and proved result
in BJ-triv, rather than a new uncached differentiability assertion.

Theorem 7.40 supplies locality on \(\{u>v\}\) for finite-energy
functions. The model proof compares vertical Cartier divisors near a
central component; the general proof passes weighted measures to the
limit using finite-energy pairing continuity. It is not just locality
for continuous weights. Proposition 7.19(v) implies that the infinite
loci of the finite-energy potentials have zero mass for all the MA
measures appearing here, so the following partitions and differences are
interpreted off a common null set. To derive comparison, put
\(w_\epsilon=\max(u,v-\epsilon)\). Positivity, locality and raw mass
give

\[V\geq\int_{\{u<v-\epsilon\}}\operatorname{MA}^{\rm raw}(v)+\int_{\{u>v-\epsilon\}}\operatorname{MA}^{\rm raw}(u),\]

and therefore

\[\int_{\{u<v-\epsilon\}}\operatorname{MA}^{\rm raw}(v)\leq\int_{\{u\leq v-\epsilon\}}\operatorname{MA}^{\rm raw}(u).\]

Letting \(\epsilon\downarrow0\) gives Theorem 7.44 in raw form,

\[\int_{\{u<v\}}\operatorname{MA}^{\rm raw}(v)\leq\int_{\{u<v\}}\operatorname{MA}^{\rm raw}(u).\]

For uniqueness, equality of MA measures gives
\(I(u,v)=\int(u-v)(\operatorname{MA}(v)-\operatorname{MA}(u))=0\). The
finite-energy integrability needed here is part of Proposition 7.19(v).
Corollary 10.4 says that \(I=0\) is equivalent to \(u-v\) being locally
constant. In its proof the energy estimates applied to \(\delta_{v}\)
and \(\delta_{v_{\rm triv}}\) show equality at every divisorial
valuation after matching the suprema; the divisorial determination
theorem then gives equality everywhere. Smooth connected \(X\) is
irreducible here, so the difference is a single constant. Multiplication
of the metric by \(e^{-c}\) corresponds to adding \(c\) to \(u\).
Comparison alone is not being substituted for this injectivity argument.

\subsection{3.5 Actual existence, with the compactness step
included}\label{actual-existence-with-the-compactness-step-included}

Let \(\nu=\mu/V\) have \(E^\vee(\nu)<\infty\). Choose a maximizing
sequence \(u_j\in\mathcal H(L)\) for \(F_\nu(u)=E(u)-\int u\,d\nu\) and
normalize \(\sup u_j=0\). The normalized functions need not remain
rational Fubini--Study functions, but lie in the continuous psh cone and
in \(\mathcal E^1\). This normalization leaves \(F_\nu\) unchanged
because \(\nu\) has mass one. Lemma 9.15 and Corollary 9.16 bound
\(J(u_j)\leq C(E^\vee(\nu)+o(1))\), where \(J(u)=\sup u-E(u)\) in the
trivial setting. For \(n\leq2\) the established envelope property gives
a weakly convergent subnet with limit \(u\in\operatorname{PSH}\).

Upper semicontinuity of \(E\) and the \(J\) bound imply
\(u\in\mathcal E^1\). Proposition 9.19 gives \(u_j\to u\) in
\(L^1(\nu)\) along this subnet. Consequently

\[F_\nu(u)\geq\limsup_j F_\nu(u_j)=E^\vee(\nu).\]

The reverse inequality is the definition of the dual supremum. Thus
\(u\) is an actual maximizing psh potential, rather than a formal
element of a completion of measures.

For completeness, the equation follows from the published uniform
derivative by an explicit argument. For PL \(f\) let
\(\delta_j=E^\vee(\nu)-F_\nu(u_j)\to0\). The obstacle-duality bound and
Theorem 8.5 give, for \(t>0\),

\[t\int f\,\operatorname{MA}(u_j)\leq t\int f\,d\nu+\delta_j+C_f t^2.\]

First let \(j\) tend to the limit, then let \(t\downarrow0\), and repeat
with \(-f\). This proves the weak convergence of
\(\operatorname{MA}(u_j)\) to \(\nu\) (Theorem 9.11). In particular the
constant maximizing sequence \(u_j=u\) proves
\(\operatorname{MA}(u)=\nu\) (Corollary 9.12). PL density determines the
measure on the full analytic space, without a support assumption.
Multiplying by \(V\) proves the raw equation. The preceding injectivity
argument completes uniqueness. This is also the proof of the applicable
direction of Theorem 12.8; Theorem 12.9 gives exactly this endpoint.

\section{4. What survives over a genuinely nontrivial
field}\label{what-survives-over-a-genuinely-nontrivial-field}

Let now \(K\) be complete, nontrivially discretely valued with residue
characteristic \(p>0\), \(X/K\) smooth projective, and \(L\) ample.
Choose a continuous semipositive reference metric \(\rho\) arising from
a positive projective model of \(L\) (a rational model of a positive
tensor power suffices). Write \(\theta=c_1(L,\rho)\) and \(V=(L^n)\). No
regular model is asserted. For nef model metrics the genuine raw measure
is

\[\sum_F b_F\deg_F(\mathcal L^{\,n})\,\delta_{v_F},\]

with the usual rational-power scaling when \(\mathcal L\) models a
tensor power. Positivity follows from nefness on each component, and
flat specialization gives the total \(V\). BE Propositions 8.3 and 8.5
and Theorems 8.4 and 8.18 supply the continuous measure and
model/Deligne-pairing calculus. Since ample \(L\) is semiample, its
continuous psh metrics are psh-regularizable by BE Theorem 7.14, as
required by Theorem 8.4. The extension uses uniform limits over an
arbitrary complete valued field. This is not the requested unbounded
operator.

For test potentials define

\[E_\theta(u)=\frac1{(n+1)V}\sum_{j=0}^n\int u\,(\theta+dd^cu)^j\wedge\theta^{n-j},\qquad J_\theta(\nu)=\sup_{u\in D_\theta}\left(E_\theta(u)-\int u\,d\nu\right).\]

Here \(D\) is the real vector space of PL/model functions containing
constants, and \(D_\theta\) consists of those with
\(\theta+dd^cu\geq0\). Thus this precise, already defined test-side raw
energy condition is \(\mu(X^{\rm an})=V\) and
\(VJ_\theta(\mu/V)<\infty\). Equality with a supremum over an eventual
unbounded \(\mathcal E^1\) cone would require a valid monotone-energy
extension and integrability passage, not just a notation change. Section
3.1 explains that passage where it is proved.

BJ-synth Sections 1.2.2 and 1.3.2 verify the synthetic data for this
projective Berkovich setting: intersection wedge measures, integration
by parts, and seminegativity by the local Hodge index theorem. In
general one must keep the Bott--Chern model quotient \(H_{\rm BC}\); an
isomorphism \(H_{\rm BC}=N^1(X)\) is stated there only under additional
hypotheses. Its positive cone is the preimage of the ample cone. This
suffices for the present ample reference and avoids assuming that
isomorphism in positive residue characteristic.

Example 2.20 derives orthogonality from BE's relative-volume theorem and
the published BGM uniform differentiability estimate. Its Definition
2.15 is

\[\lim_{u\in D_\theta,\ u<f}\int(f-u)\,\operatorname{MA}_\theta(u)=0\]

for PL \(f\), with the directed order specified in the source. It does
not first assert that a pointwise envelope is a continuous psh metric.
Theorem 2.23 then proves that any maximizing test sequence for
\(J_\theta(\nu)\) has \(\operatorname{MA}_\theta(u_j)\to\nu\)
\textbf{strongly}, and Theorem A constructs a complete quasi-metric on
the measure space. These conclusions have no restriction on
\(\operatorname{char}(K)\) or \(\operatorname{char}(\widetilde K)\) and
require no imposed resolution hypothesis in the source's projective
example. The derivative estimate yields weak convergence by the same
defect argument in Section 3.5, but without constructing a limiting psh
potential. The strong convergence of dual energies is the further source
theorem, not a consequence asserted merely from weak convergence.

Thus orthogonality at the test-energy level no longer forces retaining
BFJ-MA-v1's function-field algebraizability assumption. This does
\textbf{not} establish regularity of envelopes, compactness of the
singular psh cone, the unbounded operator, or realization of every
measure by a potential. Completing a space of measures is a different
operation from showing completeness of a particular space of usc psh
functions modulo constants. BJ-triv Theorem 12.8 proves their
relationship in its own trivial setting; it cannot be imported to
another field without its analytic inputs.

A simple genuine higher-dimensional nontrivial example helps locate the
distinction. Suppose a smooth projective \(K^\circ\)-model
\(\mathcal X\) exists with irreducible special fiber and a relatively
ample extension \(\mathcal L\) of \(L\). The special fiber has
multiplicity one. For its reference metric, the model formula gives
\(\operatorname{MA}^{\rm raw}_\rho(0)=V\delta_\xi\), where \(\xi\) is
the point attached to that component. Test energy is concave with
derivative \(\operatorname{MA}_\rho(0)=\delta_\xi\), so
\(E_\theta(u)\leq\int u\,d\delta_\xi\). The zero test gives equality in
the dual supremum; hence \(J_\theta(\delta_\xi)=0\). This solves that
reference target in any dimension with actual MA and exact mass. It is
only a particular test/continuous solution, not an existence theorem for
arbitrary finite-energy \(\mu\), nor an extension to unbounded weights.
No base change or alteration is used in this example.

\section{5. Why the recent local theories do not close this
argument}\label{why-the-recent-local-theories-do-not-close-this-argument}

Gubler--Rabinoff, arXiv:2506.13548v2, banner 17 September 2025,
introduces classically psh functions through curves and base change. Its
introduction expressly states that defining MA may require additional
hypotheses and is not the subject of the paper. Its identification of
classically psh R-PL functions with semipositivity is useful local
positivity information, but is not an approximation-independent global
\(\mathcal E^1\) equation theorem.

Chambert-Loir--Ducros, arXiv:1204.6277v2, banner 25 July 2025 and
interior date 29 July 2025, states in Section 0.2.15 that its
forms/currents theory works in every characteristic without resolution
or alterations. It develops local Bedford--Taylor and metric-curvature
products. Those characteristic-independent statements replace inputs for
that local theory. They do not by themselves identify every singular
metric in the target global psh cone with its regularizable local class,
prove global normalized compactness, or supply the required variational
existence and injectivity. No such identification is proved in this
draft.

Thuillier's thesis is potential theory on analytic curves; Berkovich's
1993 paper is \emph{Etale cohomology for non-Archimedean analytic
spaces}. Their analytic and curve-theoretic results are compatible with
positive residue characteristic. Neither is used as a replacement for
higher-dimensional envelope regularity or as an equation theorem in
higher dimension. Their PDFs were checked as sources, not invoked to
collapse the assignment to \(\mathbb P^1\).

\section{6. Proof boundaries and failed inference
chains}\label{proof-boundaries-and-failed-inference-chains}

The following passages would be invalid in the present proof and were
not made:

\begin{itemize}
\tightlist
\item
  Smooth generic fiber \(\Rightarrow\) cofinal regular SNC models in
  arbitrary positive residue characteristic. The relevant total spaces
  are one dimension larger; equal and mixed characteristic require
  distinct geometric inputs.
\item
  Resolution alone \(\Rightarrow\) the multiplier-ideal subadditivity
  and vanishing used by BFJ. Positive-characteristic test ideals require
  the hypotheses under which their generation estimates hold; the
  reported BE characteristic-\(p\) theorem imposes function-field
  descent and resolution. Mixed characteristic is not covered by that
  report.
\item
  Alterations or a universal homeomorphism \(\Rightarrow\) a descended
  finite-energy operator. Even a topological statement does not prove
  preservation of the ample class, psh cone, masses, energy, derivative,
  unbounded products or comparison.
\item
  Strong convergence of genuine model MA measures \(\Rightarrow\) a psh
  solution of their limit equation. A realization/completeness theorem
  for actual potentials is missing.
\item
  Local current products \(\Rightarrow\) global variational existence
  for every finite-dual-energy measure. Compactness, envelope
  compatibility and unbounded injectivity remain separate.
\end{itemize}

These are logical gaps in attempted extensions, not mathematical
counterexamples to the target. No nonexistence assertion is proved.

\section{7. Concrete remaining original
obligations}\label{concrete-remaining-original-obligations}

For complete nontrivially discretely valued \(K\) with
\(\operatorname{char}(\widetilde K)=p>0\), smooth connected projective
\(X/K\) in arbitrary dimension, ample \(L\), and all positive Radon
\(\mu\) of mass \(V\) and finite dual energy, the following work is
still required:

\begin{enumerate}
\def\labelenumi{\arabic{enumi}.}
\tightlist
\item
  \textbf{Singular psh theory and compactness.} Supply a definition
  equivalent to the intended model-semipositive theory, monotone
  approximation, divisorial determination where used, and compactness
  modulo constants. Replace the SNC-cofinality, retraction and toroidal
  equicontinuity arguments, or prove an alternative with all hypotheses.
  The original residue-characteristic-zero theorems cannot be cited on
  the changed field.
\item
  \textbf{Model positivity and envelopes.} Prove the required
  positivity/nef approximation and continuity or an adequate variational
  substitute for envelopes without an unproved resolution hypothesis or
  extra function-field descent. Record any restriction if such an input
  is only conditional. Cover mixed characteristic separately from equal
  characteristic \(p\).
\item
  \textbf{Unbounded operator.} Construct the genuine positive mixed MA
  products on \(\mathcal E^1\); prove finiteness/integrability, raw full
  mass \(V\), independence of decreasing approximants or truncations,
  and compatibility with model and continuous MA. Test-side measure
  completion does not meet this item.
\item
  \textbf{Exact energy and derivative on that space.} Establish the
  monotone energy extension, equality of the test and \(\mathcal E^1\)
  dual suprema, and the energy/obstacle derivative at finite-energy
  potentials. The modern test orthogonality result is an available
  input, not by itself this passage.
\item
  \textbf{Existence of an actual potential.} Show that a normalized
  maximizing family has a limit in the stated psh cone of finite energy,
  that the measure integral passes to the limit, and that the derivative
  characterizes its genuine MA as the full \(\mu\). Do this for
  arbitrary support without restricting and renormalizing the measure.
\item
  \textbf{Comparison and uniqueness.} Prove locality/comparison for
  unbounded finite-energy weights and injectivity modulo constants using
  a valid energy or domination argument. Continuous-metric uniqueness
  does not establish finite-energy uniqueness.
\end{enumerate}

No new indispensable passage meeting all six items over the target
fields has been proved here. The unconditional theorem in Section 3 is
narrower and remains partial. Even the unresolved higher-dimensional
trivially valued envelope property is a separate limitation, not a
replacement for these nontrivial-field obligations.

\section{8. Reproduction and evidence}\label{reproduction-and-evidence}

The original JSON is preserved byte for byte in
\texttt{input-original.json}; the decoded assignment and task are in
\texttt{assignment-original.txt} and \texttt{task-original.json}.
\texttt{source\_audit.py} reads only the supplied input and the ten
listed public caches, checks all 20 PDF/text hashes and six abstract
hashes, copies the sources to \texttt{source-extracts}, regenerates PDF
text with the installed \texttt{pdftotext}, and saves 52 literal
extracts with the published wording and precise line/page locations in
\texttt{source-extracts/source-provenance.json}. The original cache
hashes, version banners, application checks, and extract texts are
present in that JSON. OCR typos are retained rather than silently
corrected in literal extracts.

The PDF page images \texttt{source-extracts/BJ22-page-62.png},
\texttt{BJ22-page-97.png} and \texttt{BJ22-page-122.png} were inspected
for the model-dimension ambiguity, the energy derivative, and the
existence theorem. Their visual checks and hashes are recorded in
provenance. In particular, page 62 says \(\dim X\leq2\) and
\(\dim\mathcal X\leq3\); page 122 states the unconditional
positive-characteristic dimension bound.

Rerun \texttt{python3\ source\_audit.py} and
\texttt{python3\ build\_and\_check.py} in this workspace. The latter
generates a standalone \texttt{P16.3-paper.tex} from this argument using
the existing Pandoc, runs the existing \texttt{pdflatex} twice
sequentially, saves each actual stdout/stderr diagnostic stream and TeX
log, and checks the input preservation and exact status schema.
\texttt{status.json} has exactly
\texttt{id,status,exact\_result,proof\_file,concrete\_gaps}. Source
hashes and compilation verify source identity and artifact execution
only; they do not verify the mathematical arguments. No heavy
computation, installation, network fetch, model call, worker, or
external message is used.

\section{References (exact supplied
artifacts)}\label{references-exact-supplied-artifacts}

\begin{enumerate}
\def\labelenumi{\arabic{enumi}.}
\tightlist
\item
  BFJ-PSH: arXiv:1201.0187v2, \emph{Singular semipositive metrics in
  non-Archimedean geometry}.
\item
  BFJ-MA-v1: arXiv:1201.0188v1, \emph{Solution to a non-Archimedean
  Monge-Ampere equation}.
\item
  BE: arXiv:1805.01016v3, \emph{Spaces of norms, determinant of
  cohomology and Fekete points in non-Archimedean geometry}.
\item
  BJ-triv: arXiv:1801.08229v4, \emph{Global pluripotential theory over a
  trivially valued field}.
\item
  BJ-synth: arXiv:2307.01697v1, \emph{Measures of finite energy in
  pluripotential theory: a synthetic approach}.
\item
  PS-hybrid: arXiv:2209.04879v2, \emph{Global pluripotential theory on
  hybrid spaces}.
\item
  Gubler--Rabinoff: arXiv:2506.13548v2, \emph{Classically psh and
  pluriharmonic functions on Berkovich spaces}.
\item
  Thuillier: tel-00010990v1, \emph{Theorie du potentiel sur les courbes
  en geometrie analytique non archimedienne. Applications a la theorie
  d'Arakelov}, thesis, 2005.
\item
  Berkovich: \emph{Etale cohomology for non-Archimedean analytic
  spaces}, Publications mathematiques de l'IHES 78 (1993), 5--161,
  supplied Numdam PDF.
\item
  Chambert-Loir--Ducros: arXiv:1204.6277v2, \emph{Formes differentielles
  reelles et courants sur les espaces de Berkovich}.
\end{enumerate}

\end{document}
